% =====================================================================
%  Chương 2 — Mục 2 của đề: các mô hình phát triển của CNN
%  Mã nguồn: a05/blocks.py · Notebook: notebooks/01_kien_truc_phat_trien.ipynb
% =====================================================================
\chapter{Các mô hình phát triển của CNN}
\label{ch:kientruc}

Chương 1 cho thấy CNN là hợp của các hàm. Từ năm 1998 đến nay, mỗi kiến trúc mới thường chỉ đổi \emph{một}
trong các hàm ấy, hoặc đổi cách chúng được nối với nhau, để sửa một điểm yếu cụ thể của kiến trúc trước.
Chương này đi qua chín mốc chính (Hình~\ref{fig:dongthoigian}). Với mỗi mốc, chương nêu vấn đề nó giải quyết,
viết cơ chế mới thành công thức, rồi đưa ra một lớp PyTorch ngắn trong \texttt{a05/blocks.py}. Notebook 01
chạy mỗi lớp một lượt trên tensor giả để kiểm tra shape và đếm tham số; mục này không huấn luyện.

\begin{figure}[H]
  \centering
  \hinh{f2_dong_thoi_gian}
  \caption{Dòng thời gian các kiến trúc CNN và cơ chế mà mỗi kiến trúc đưa vào}
  \label{fig:dongthoigian}
\end{figure}

\section{LeNet-5 (1998): khuôn mẫu Conv, Pool, FC}

LeCun và cộng sự~\cite{lecun1998} đặt ra khuôn mẫu mà mọi CNN sau này vẫn giữ: vài lần lặp
(tích chập, hàm kích hoạt, pooling) để trích đặc trưng, rồi các tầng nối đầy đủ để phân loại,
\[
  f = \mathrm{FC}_3\circ\tanh\circ\mathrm{FC}_2\circ\tanh\circ\mathrm{FC}_1\circ\mathrm{flatten}\circ
      (\mathrm{pool}\circ\tanh\circ\mathrm{conv}_{5\times5})^{2}.
\]
Bản dựng lại cho ảnh $32\times32\times3$ có \SL{blocks/blocks/lenet/params} tham số.

\maNguon{blocks__LeNet5}{LeNet-5}

\section{AlexNet (2012): sâu hơn, ReLU và Dropout}

LeNet dùng $\tanh$, có đạo hàm nhỏ ở hai đầu. Khi mạng sâu thêm, tích của nhiều đạo hàm nhỏ làm gradient
ở các tầng đầu gần như bằng 0. AlexNet~\cite{krizhevsky2012} thay $\tanh$ bằng $\mathrm{ReLU}(z)=\max(0,z)$,
có đạo hàm bằng 1 trên miền dương. Để chống quá khớp ở các tầng nối đầy đủ lớn, AlexNet dùng Dropout:
khi huấn luyện, mỗi đơn vị bị đặt về 0 với xác suất $p$, tức nhân đầu ra với mặt nạ ngẫu nhiên
$\mathbf{m}\sim\mathrm{Bernoulli}(1-p)$. Bản thu nhỏ trong \texttt{blocks.py} có \SL{blocks/blocks/alexnet/params}
tham số, phần lớn nằm ở tầng nối đầy đủ đầu tiên.

\maNguon{blocks__AlexNetMini}{AlexNet thu nhỏ cho ảnh $32\times32$}

\section{VGG (2014): chồng các tích chập $3\times3$}

Simonyan và Zisserman~\cite{simonyan2015} chỉ dùng một loại bộ lọc: $3\times3$. Hai tầng $3\times3$ liên tiếp
có vùng nhìn $5\times5$, giống một tầng $5\times5$, nhưng với $C$ kênh vào và ra thì cần $2\cdot9C^2=18C^2$
trọng số thay vì $25C^2$, đồng thời có thêm một hàm phi tuyến ở giữa. Với $C=64$, notebook 01 đếm được
\SL{blocks/compare/conv5} trọng số cho một conv $5\times5$ và \SL{blocks/compare/two3} cho hai conv $3\times3$,
ít hơn \SL{blocks/compare/vgg_saving}\,\%. Một khối VGG là hàm
$\mathrm{pool}\circ(\mathrm{ReLU}\circ\mathrm{conv}_{3\times3})^{n}$.

\maNguon{blocks__VGGBlock}{Khối VGG}

\section{Inception (2014): nhiều cỡ bộ lọc song song}

Chọn bộ lọc $3\times3$ hay $5\times5$ là chọn trước một cỡ vùng nhìn. Inception~\cite{szegedy2015} không chọn:
nó chạy song song bốn nhánh rồi nối kết quả theo chiều kênh,
\[
  y = \mathrm{concat}\bigl(f_{1\times1}(x),\; f_{3\times3}(x),\; f_{5\times5}(x),\; f_{\mathrm{pool}}(x)\bigr).
\]
Để nhánh $3\times3$ và $5\times5$ không quá đắt, một conv $1\times1$ giảm số kênh trước đó. Conv $1\times1$ là một
tầng nối đầy đủ áp dụng riêng tại từng điểm ảnh, trộn các kênh mà không nhìn sang điểm bên cạnh.

\maNguon{blocks__InceptionBlock}{Khối Inception}

\section{ResNet (2015): đường tắt}

Khi VGG được làm sâu hơn khoảng 20 tầng, sai số trên cả tập \emph{huấn luyện} lại tăng lên. Đây không phải
quá khớp mà là mạng khó tối ưu. He và cộng sự~\cite{he2016} đổi mục tiêu của mỗi khối: thay vì học trực tiếp
ánh xạ $H(x)$, khối học phần dư $F(x) = H(x) - x$,
\[
  y = \mathrm{ReLU}\bigl(F(x) + x\bigr).
\]
Hai hệ quả trực tiếp của công thức này. Thứ nhất, nếu $F\equiv0$ thì khối là ánh xạ đồng nhất, nên thêm một khối
không thể làm mạng tệ đi. Notebook 01 kiểm tra điều này bằng cách đặt $\gamma$ của BatchNorm cuối về 0:
đầu ra khớp đầu vào (\SL{blocks/identity_ok}). Thứ hai, đạo hàm
$\partial y/\partial x = \partial F/\partial x + I$ luôn có thành phần $I$, nên gradient đi thẳng qua đường tắt
mà không bị thu nhỏ qua từng tầng. Khối minh hoạ với 64 kênh có \SL{blocks/blocks/residual/params} tham số.

\maNguon{blocks__ResidualBlock}{Khối Residual}

\section{DenseNet (2017): nối mọi tầng với nhau}

DenseNet~\cite{huang2017} đẩy ý tưởng đường tắt đi xa hơn. Mỗi tầng nhận \emph{tất cả} đầu ra trước nó, nối
theo kênh chứ không cộng:
\[
  x_l = H_l\bigl([x_0, x_1, \dots, x_{l-1}]\bigr).
\]
Mỗi tầng chỉ thêm $k$ kênh mới (tốc độ tăng trưởng, \emph{growth rate}), nên các tầng có thể rất hẹp. Khối trong
notebook bắt đầu từ 16 kênh, qua 4 tầng với $k=12$, ra $16+4\cdot12=64$ kênh, dùng
\SL{blocks/blocks/dense/params} tham số.

\maNguon{blocks__DenseBlock}{Khối Dense}

\section{MobileNet (2017): tách tích chập thành hai bước}

Một tầng conv $3\times3$ thường vừa lọc theo không gian vừa trộn kênh, tốn $9\,C_{\mathrm{in}}C_{\mathrm{out}}$
trọng số. MobileNet~\cite{howard2017} tách hai việc đó: \emph{depthwise} lọc từng kênh riêng
($9\,C_{\mathrm{in}}$ trọng số), rồi \emph{pointwise} $1\times1$ trộn kênh ($C_{\mathrm{in}}C_{\mathrm{out}}$).
Với $64\to128$ kênh, conv thường cần \SL{blocks/compare/std} trọng số, còn bản tách chỉ cần
\SL{blocks/compare/dws}, ít hơn \SL{blocks/compare/dws_ratio} lần.

\maNguon{blocks__DepthwiseSeparable}{Tích chập depthwise separable}

\section{SE-Net (2017) và CBAM (2018): chú ý}

Các kiến trúc trên đều coi mọi kênh đặc trưng như nhau. Squeeze-and-Excitation~\cite{hu2018} cho mạng tự
học một trọng số cho từng kênh, tính từ chính dữ liệu đầu vào:
\[
  \mathbf{s} = \mathrm{GAP}(x)\in\mathbb{R}^{C},\qquad
  \mathbf{w} = \sigma\bigl(W_2\,\mathrm{ReLU}(W_1\mathbf{s})\bigr)\in(0,1)^{C},\qquad
  y_c = w_c\, x_c .
\]
Bước \emph{squeeze} nén mỗi kênh về một số bằng trung bình toàn cục, bước \emph{excitation} là một MLP nhỏ với
tỉ lệ nén $r$ ($W_1\in\mathbb{R}^{C/r\times C}$). Khối này rất rẻ: với 64 kênh chỉ thêm
\SL{blocks/blocks/se/params} tham số. Trên tensor ngẫu nhiên, notebook 01 đo được các hệ số nằm trong
khoảng (\SL{blocks/se_min}; \SL{blocks/se_max}). CBAM~\cite{woo2018} làm thêm một bước chú ý theo
\emph{không gian}: gộp các kênh bằng trung bình và max, qua một conv $7\times7$ và sigmoid để có trọng số cho từng vị trí.

\maNguon{blocks__SEBlock}{Khối Squeeze-and-Excitation}
\maNguon{blocks__CBAM}{CBAM: chú ý theo kênh rồi theo không gian}

\section{Vision Transformer (2020): bỏ tích chập}

ViT~\cite{dosovitskiy2021} cắt ảnh thành các ô $p\times p$, chiếu mỗi ô thành một vectơ (\emph{token}), rồi dùng
self-attention để mọi ô nhìn thấy mọi ô khác ngay từ tầng đầu. Tích chập chỉ nhìn một vùng lân cận nhỏ; self-attention thì
tính trọng số giữa hai ô bất kỳ:
\[
  \mathrm{Attn}(Q,K,V) = \mathrm{softmax}\!\left(\frac{QK^{\top}}{\sqrt{d}}\right)V .
\]
Điều thú vị là bước nhúng ô chính là một tầng conv có kích thước bộ lọc bằng bước trượt bằng $p$
(Mã nguồn~\ref{code:blocks__PatchEmbed}). Với $p=4$, ảnh $32\times32$ thành $8\cdot8=64$ token.

\maNguon{blocks__PatchEmbed}{Nhúng ô ảnh của ViT là một tầng tích chập}
\maNguon{blocks__TinyViTBlock}{Một khối encoder của ViT}

\section{Tổng hợp}

Hình~\ref{fig:coche} vẽ sáu cơ chế chính dưới dạng sơ đồ khối, và Bảng~\ref{tab:khoi} ghi shape cùng số tham số
mà notebook 01 đo được cho từng khối.

\begin{figure}[H]
  \centering
  \hinh{f2_co_che}
  \caption{Sáu cơ chế kiến trúc viết thành sơ đồ khối}
  \label{fig:coche}
\end{figure}

\begin{table}[H]
  \centering
  \caption{Shape và số tham số của các khối minh hoạ (notebook 01, tensor giả, lô 2)}
  \label{tab:khoi}
  \input{generated/tab_blocks}
\end{table}

\begin{phathien}{Ba mô hình phát triển được chọn cho Chương 4}
Trong chín mốc trên, ba cơ chế được đem đi huấn luyện ở Chương~\ref{ch:mohinh}: chồng conv $3\times3$ kèm
BatchNorm (VGG), đường tắt (ResNet) và chú ý theo kênh (SE). Lý do chọn ba cơ chế này: mỗi cơ chế thêm được
vào mô hình trước mà không phải đổi phần còn lại, nên mỗi bước so sánh chỉ khác đúng một cơ chế.
Inception, DenseNet và ViT đổi cả khung kiến trúc, nên không so sánh theo cách này được.
\end{phathien}
